% Nouveau EGA V (cote 25): archive cote 25, site page 92 (PDF page 93).
% typescript. Draft compared: none.
% AI-transcribed by Claude Opus 5.5 (claude-opus-5-5, adaptive thinking, effort high), Message Batches API, batch ega25b (msgbatch_01A7n5YuXubUPNK1zsepi8RG), 2026-11-13; not yet checked by eye.
\scanpage{25/92}
\mspage{86}

tout au moins sont valables sous d es conditions plus générales que \inserted{pour} la
famille des hyperplans (ou des hypersurfaces de degré donné) de l'espace
projectif. Il me semblerait opportun d'adopter un tel point de vue
axiomatique d\over{an}{\unsure{è}}s les premiers n$^{\text{o}}$s de la rédaction en forme. Je n'ai pas
tiré au clair pour le moment s'il y a bien moyen de donner une généralisa%
tion dans ce sens de Bertini-Zariski (donc des résultats des n$^{\text{o}}$s 4 et 6),
et j'ai écrit aux compétences (Serre, Zariski) pour leur demander s'ils
avaient connaissance d'une telle extension; j'ai bien l'impression en
effet que des hypothèses de nature différentielle simple, du genre de cel
% NOTE: the end of the line runs off the right edge of the scan after "cel"
les envisagées ci-dessous, doivent suffire à impliquer Bertini-Zariski.
Si les compétences ne peuvent nous informer de façon satisfaisante, il
faudrait essayer de tirer la chose au clair par nos propres moyens.

On part d'un diagramme \inserted{commutatif} de morphismes de présentation finie

\margin{(D)}
\diagram{
\begin{tikzcd}
P & \underline{P} \arrow[l] \arrow[d] \\
S \arrow[u, dash] & G \arrow[l]
\end{tikzcd}
}

(dans le cas d'application principal, P est un fibré projectif, G une
grasm\unsure{a}nienne dudit, $\underline{P}$ le préschéma d'incidence). Dans lescas lesplus
importants, le morphisme correspondant $\underline{P} \longrightarrow P\times_S G$ sera une immersion
fermée, et on considère G comme \xout{définissant une} \inserted{schéma des paramètres d'une} familled e sous-
préschémas fermés des fibres de P sur S, d e façon précise si \hand{$\zeta$} \hand{$\in$} G,
$\underline{P}_{\hand{\zeta}}$ est un sous-préschéma fermé de $P_{s\ k(\hand{\zeta})}$ , où s est le point de S en
dessous de \hand{$\zeta$} .(D'ailleurs, pour la plupart des énoncés dans ce contexte,
on aura sans doute S=Spec(k)\hand{)}. Dans le cas général, on peut encore con-%
sidérer G comme schéma des paramètres d'une famille de préschémas
au dessus des fibres de P sur S, à \hand{$\zeta$} correspondant $\underline{P}_{\hand{\zeta}}$ au dessus de
$P_{s\ k(\hand{\zeta})}$ . Bien entendu, au lieu de prendre pour \hand{$\zeta$} un point \xout{de} (absolu)

\margin{92}
